\documentclass[11pt]{article}

% cs250preamble.tex --- shared preamble for CS 250 course notes
%
% This file is \input by main.tex and by each individual module
% file so that each module can still compile standalone. Any
% package or custom-environment change should be made here, not
% in the individual module files.

\usepackage[T1]{fontenc}
\usepackage[margin=1in]{geometry}
\usepackage{amsmath, amssymb, amsthm}
\usepackage{enumitem}
\usepackage{hyperref}
\usepackage{booktabs}
\usepackage{listings}
\usepackage{xcolor}
\usepackage{tikz}
\usetikzlibrary{shapes.geometric, shapes.gates.logic.US, shapes.gates.logic.IEC, arrows.meta, positioning, calc, trees}
\usepackage{bussproofs}
\usepackage{mdframed}

\lstset{
  basicstyle=\ttfamily\small,
  frame=single,
  breaklines=true,
  columns=fullflexible,
  mathescape=true
}

\theoremstyle{definition}
\newtheorem{theorem}{Theorem}[section]
\newtheorem{definition}[theorem]{Definition}
\newtheorem{example}[theorem]{Example}

\hypersetup{
  colorlinks=true,
  linkcolor=blue!60!black,
  urlcolor=blue!60!black
}

% Key-takeaway box: used at the end of major sections and at the end of
% each module to summarise the main points. Expects \item bullets inside.
\newmdenv[
  linewidth=0.8pt,
  linecolor=blue!40!black,
  backgroundcolor=blue!3,
  skipabove=8pt,
  skipbelow=8pt,
  innertopmargin=7pt,
  innerbottommargin=7pt,
  innerleftmargin=9pt,
  innerrightmargin=9pt
]{keytakeawaybox}
\newenvironment{keytakeaway}{%
  \begin{keytakeawaybox}%
  \noindent\textbf{Key takeaways.}%
  \begin{itemize}[leftmargin=1.3em, topsep=2pt, itemsep=1pt, label=\textbullet]%
}{%
  \end{itemize}%
  \end{keytakeawaybox}%
}

% Summary/looking-ahead environment: used once at the end of each module
% to wrap up what the module built and point forward.
\newmdenv[
  linewidth=0.8pt,
  linecolor=gray!60,
  backgroundcolor=gray!5,
  skipabove=8pt,
  skipbelow=8pt,
  innertopmargin=7pt,
  innerbottommargin=7pt,
  innerleftmargin=9pt,
  innerrightmargin=9pt
]{summarybox}

% Sidebar environment: used for tangential asides---historical notes,
% pointers to other areas of math/CS, cross-paradigm remarks---that
% enrich the main narrative without being essential to it. Takes one
% argument: the sidebar's title.
\newmdenv[
  linewidth=0.6pt,
  linecolor=gray!60,
  backgroundcolor=gray!4,
  skipabove=8pt,
  skipbelow=8pt,
  innertopmargin=6pt,
  innerbottommargin=6pt,
  innerleftmargin=9pt,
  innerrightmargin=9pt
]{sidebarbox}
\newenvironment{sidebar}[1]{%
  \begin{sidebarbox}%
  \noindent\textbf{Sidebar: #1.}\par\medskip%
}{%
  \end{sidebarbox}%
}

% Reading-map environment: used at the top of each numbered module to
% indicate Epp coverage, prerequisites, relevant intermezzi, and
% suggested practice problems. Expects four \rmentry{label}{body}
% items inside (or arbitrary paragraphs).
\newmdenv[
  linewidth=0.8pt,
  linecolor=black!55,
  backgroundcolor=yellow!4,
  skipabove=10pt,
  skipbelow=14pt,
  innertopmargin=9pt,
  innerbottommargin=9pt,
  innerleftmargin=11pt,
  innerrightmargin=11pt
]{readingmapbox}
\newcommand{\rmentry}[2]{\par\noindent\textbf{#1}\hspace{0.4em}#2\par}
\newenvironment{readingmap}{%
  \begin{readingmapbox}%
  \noindent{\bfseries\large Reading map}%
  \vspace{2pt}%
  \setlength{\parskip}{4pt plus 1pt}%
}{%
  \end{readingmapbox}%
}


\title{CS 250 --- Intermezzo: Quantifiers as Games}
\author{}
\date{}

\begin{document}
\maketitle

In Module 4 we defined quantifiers two ways: as functions into Bool, and as Big And / Big Or over a set. Both of those descriptions are perfectly correct, but they're also a bit static---they tell you what quantifiers \emph{mean} but not what it \emph{feels like} to work with them. There's a third perspective that is more dynamic, more intuitive, and connects beautifully to both programming and to modern cryptography. It's the idea that logic is a \emph{game}.

\section{The game interpretation}

Imagine two players sitting across from each other. One is the \textbf{Prover} (that's you---the person trying to show the statement is true). The other is the \textbf{Challenger} (think of them as Nature, or an Opponent, or a very skeptical friend). The logical formula being evaluated determines the rules of the game. Prover wins if, at the end of the game, the resulting atomic statement is true. Challenger wins if it's false.

Here's how each connective translates into a game move:

\begin{itemize}
  \item $\exists x.\; P(x)$: \textbf{Prover's move.} Prover picks a witness $x$, and then the game continues with $P(x)$. Prover wins if $P(x)$ turns out to be true for the chosen $x$. This is the ``I can find one'' quantifier---Prover gets to choose.

  \item $\forall x.\; P(x)$: \textbf{Challenger's move.} Challenger picks \emph{any} $x$ they want---the nastiest, most adversarial $x$ they can think of---and then the game continues with $P(x)$. Prover wins if $P(x)$ is true \emph{regardless} of what Challenger picked. This is the ``it works for everything'' quantifier---Challenger gets to choose, and Prover has to survive.

  \item $A \wedge B$: \textbf{Challenger's move.} Challenger picks which of $A$ or $B$ to verify. Prover must be able to win \emph{both} games, so it doesn't matter which one Challenger chooses. Think of it as: ``I claim both---go ahead, quiz me on either one.''

  \item $A \vee B$: \textbf{Prover's move.} Prover picks which of $A$ or $B$ to demonstrate. Only one needs to hold, and Prover gets to pick the one they can actually win. Think of it as: ``I claim at least one of these---let me show you which.''

  \item $A \to B$: \textbf{Challenger provides, Prover responds.} Challenger hands over evidence for $A$, and Prover must use it to produce evidence for $B$. If you're a programmer, this should look familiar: it's a function. Challenger provides the input, Prover computes the output.

  \item $\lnot A$: \textbf{Role swap.} The players switch seats. Prover becomes Challenger, Challenger becomes Prover, and they play the game for $A$ with reversed roles. Negation flips who's trying to do what.
\end{itemize}

This might seem like just a cute metaphor, but it's not---it's a genuine mathematical framework called \emph{game semantics}\footnote{Game semantics was developed seriously in the 1990s by Abramsky, Jagadeesan, Malacaria, Hyland, Ong, and others. It provides a fully compositional semantics for programming languages and logics---not just a teaching analogy.} and it gives us the exact same truth values as the set-based semantics from Module 4. A formula is true precisely when Prover has a \emph{winning strategy}---a way to play that wins no matter what Challenger does.

\section{Nested quantifiers as alternating turns}

The game interpretation really shines when we look at nested quantifiers, because nested quantifiers correspond to a game with \emph{multiple rounds}---the players alternate turns, and the whole formula describes a back-and-forth.

Let's revisit the example from Module 4 where we saw that quantifier order matters. Let $P(x, y) = $ ``$y > x$'' over $\mathbb{R}$.

\medskip
\noindent\textbf{Game 1:} $\forall x : \mathbb{R}.\; \exists y : \mathbb{R}.\; y > x$

\begin{enumerate}
  \item Challenger picks any real number $x$. (Challenger's move---this is $\forall$.)
  \item Prover sees what $x$ was chosen and picks $y$ in response. (Prover's move---this is $\exists$.)
  \item Prover wins if $y > x$.
\end{enumerate}

Can Prover always win? Yes! Prover's strategy is simple: whatever $x$ Challenger picks, respond with $y = x + 1$. Since $x + 1 > x$ for every real number, this strategy works against every possible challenge. The formula is \textbf{true}.

The key thing to notice: Prover's choice of $y$ gets to \emph{depend on} $x$. Prover moves second, so they have information. Their winning strategy is a \emph{function} from challenges to responses: $f(x) = x + 1$.

\medskip
\noindent\textbf{Game 2:} $\exists y : \mathbb{R}.\; \forall x : \mathbb{R}.\; y > x$

\begin{enumerate}
  \item Prover picks a real number $y$. (Prover's move---this is $\exists$.)
  \item Challenger sees what $y$ was chosen and picks $x$ in response. (Challenger's move---this is $\forall$.)
  \item Prover wins if $y > x$.
\end{enumerate}

Can Prover win? No. Whatever $y$ Prover commits to, Challenger can just pick $x = y + 1$, and then $y > y + 1$ is false. Prover has to commit \emph{before} seeing the challenge, and no single constant $y$ can be bigger than every real number. The formula is \textbf{false}.

\medskip

See the difference? In Game 1, Prover has a \emph{strategy}---a function that responds to each challenge. In Game 2, Prover needs a \emph{constant}---a single value that works universally. Strategies are much more powerful than constants, which is why $\forall x.\;\exists y$ is a weaker claim than $\exists y.\;\forall x$ (the second implies the first, but not vice versa).

If you're a programmer, the distinction is really clean:
\begin{itemize}
  \item $\forall x.\;\exists y.\; y > x$ is like writing a function:
\begin{lstlisting}
def respond(x):
    return x + 1    # always bigger than the input
\end{lstlisting}
  You get to see the input before choosing your output. Easy.

  \item $\exists y.\;\forall x.\; y > x$ is like declaring a global constant:
\begin{lstlisting}
y = ???    # must be bigger than ALL possible inputs
\end{lstlisting}
  You have to commit before seeing any input. Impossible (for the reals).
\end{itemize}

This is the same intuition we built in Module 4, but the game framing makes the ``why'' visceral: it's about \emph{who moves first}. The order of quantifiers determines the order of play, and the order of play determines how much information each player has when they make their choice.

\section{Games and constructive logic}

There's a philosophical point lurking here that connects back to something we touched on in Module 5: the distinction between constructive and non-constructive proofs.

The game interpretation is inherently \emph{constructive}. To win a game, Prover can't just argue abstractly that a winning move exists somewhere---Prover has to actually \emph{make the move}. When the formula says $\exists x.\; P(x)$, Prover has to put a concrete $x$ on the table. When it says $A \vee B$, Prover has to point to one of $A$ or $B$ and win that game. No hand-waving allowed.

This is exactly the spirit of constructive mathematics: to prove something exists, you must build it. To prove a disjunction, you must know which side is true.

Classical logic, by contrast, gives Prover an extra power: \emph{strategy-stealing arguments}. In classical logic, Prover is allowed to reason like this: ``Suppose I can't win. Then Challenger must have a winning strategy. But if Challenger has a winning strategy, then [some contradiction arises], so actually I \emph{can} win.'' The conclusion is that Prover wins, but the argument never actually described a move for Prover to make. It proved existence without providing a witness.

This connects directly to the law of excluded middle, $p \vee \lnot p$. In game terms, this says: ``Either Prover can demonstrate $p$, or Prover can demonstrate $\lnot p$.'' Classically, this is just assumed---one of the two must be the case. But constructively, Prover has to actually play one of the two games and win it. For many propositions (especially undecidable ones from computability theory) there's no general way to know which side to pick, so $p \vee \lnot p$ isn't something you get for free.

The link to the Curry--Howard intermezzo is worth flagging explicitly: Prover's winning strategy \emph{is} a program, in essentially the same sense that Curry--Howard's ``a proof is a program'' is. A strategy for $\forall x.\,\exists y.\, P(x, y)$ is a function that turns each challenge $x$ into a response $y$---which is exactly what Curry--Howard says a proof of that formula is. Game semantics and Curry--Howard are two faces of the same computational reading of logic, and students who internalize one usually find the other suddenly clicks.

If you remember the discussion of constructive existence proofs from Module 5---where we said that a constructive proof gives you an algorithm while a non-constructive proof gives you a theorem but no code---this is the same idea wearing different clothes. Game semantics makes it vivid: constructive logic is the version of the game where Prover actually has to play. Classical logic is the version where Prover is sometimes allowed to win by saying ``if I couldn't win, that would be absurd.''\footnote{If that sounds like cheating, welcome to the intuitionist club. If it sounds perfectly reasonable, welcome to the classical club. Both clubs have cookies; they just disagree about whether the cookies have been constructed or merely shown to exist.}

\section{Why this matters}

You might be wondering: is this just a nice way to think about quantifiers, or does it actually \emph{go} somewhere? It goes a lot of places.

\textbf{Interactive proof systems.} In theoretical computer science and cryptography, the Prover-vs-Verifier model is the literal foundation of \emph{interactive proof systems}. A Prover tries to convince a Verifier that a statement is true, and the Verifier is allowed to ask questions (challenges). The statement is ``provable'' if there's a Prover strategy that convinces the Verifier, and ``sound'' if no cheating Prover can fool the Verifier into accepting a false statement. This is essentially the same two-player game we've been discussing, just taken very seriously.

\textbf{Zero-knowledge proofs.} An especially remarkable variant: what if Prover wants to convince Verifier that a statement is true \emph{without revealing any information about why it's true}? This is a zero-knowledge proof, and it's the cryptographic technology behind privacy-preserving blockchain protocols, anonymous credentials, and a growing chunk of modern cryptography. The game-theoretic structure is essential---the ``zero knowledge'' property is defined in terms of what Verifier can learn from the interaction.

\textbf{Complexity theory.} One of the landmark results in theoretical computer science is that $\mathsf{IP} = \mathsf{PSPACE}$: the class of problems solvable by an interactive proof system (where a polynomial-time Verifier interacts with an all-powerful Prover) is exactly the class of problems solvable with polynomial space. This was a shocking result when it was proved in 1992, and the entire framework rests on formalizing computation as a game between Prover and Verifier.

The point is that the two-player game isn't just a pedagogical metaphor. It's a lens through which serious mathematics, computer science, and cryptography are actually done. When you see quantifiers as games, you're not just learning a trick for reading formulas---you're learning the conceptual vocabulary of a significant part of modern theory.

\section*{Exercises}

\bigskip
\textbf{Exercise 1.} \label{ex:games:1} Consider the statement
\[
\forall \epsilon > 0.\; \exists \delta > 0.\; \forall x.\; (|x - 3| < \delta \;\to\; |x^2 - 9| < \epsilon).
\]
This is the epsilon-delta definition of continuity for $f(x) = x^2$ at $x = 3$.

Describe the game: who moves first? What does each player choose at each turn? What must Prover achieve in order to win? (You do not need to find Prover's actual strategy---just describe the structure of the game.)

\bigskip
\textbf{Exercise 2.} \label{ex:games:2} In the game for $\exists x.\; \forall y.\; (x + y = y)$, what is Prover's winning move? (Hint: Prover moves first. What value of $x$ makes the inner statement true no matter what Challenger picks for $y$?)

\bigskip
\textbf{Exercise 3.} \label{ex:games:3} Consider the formula $\forall x \in \mathbb{N}.\; \exists y \in \mathbb{N}.\; y > x$ over the natural numbers. Describe the game in detail (who moves first, who moves second, what they choose, what Prover must achieve). Then give Prover's winning strategy as an explicit function $f$ such that $f(x)$ is Prover's response when Challenger picks $x$.

\bigskip
\textbf{Exercise 4.} \label{ex:games:4} Now consider $\exists y \in \mathbb{N}.\; \forall x \in \mathbb{N}.\; y > x$. Describe the game and explain, in game-theoretic terms, exactly why Prover \emph{cannot} win, no matter what value of $y$ they pick first. Connect your answer to the difference between a strategy and a constant in the discussion above.

\bigskip
\textbf{Exercise 5.} \label{ex:games:5} (Quantifier games for arithmetic identities.) Describe the game corresponding to the statement $\forall n \in \mathbb{N}.\; \forall m \in \mathbb{N}.\; n + m = m + n$. Both quantifiers are universal, so Challenger gets to pick \emph{both} values. What does Prover have to demonstrate? In what sense is this game ``trivial,'' even though there are infinitely many possible challenges?

\bigskip
\textbf{Exercise 6.} \label{ex:games:6} (Strategies as functions.) The Skolem form of $\forall x.\; \exists y.\; R(x, y)$ replaces the existential by a function: $\exists f.\; \forall x.\; R(x, f(x))$. Explain how this transformation matches what the game interpretation says about Prover's strategy. (Hint: in the game for $\forall x.\; \exists y.\; R(x, y)$, Prover responds to each $x$ with some $y$. Bundle all those responses into a single object. What is that object?)

\bigskip
\textbf{Exercise 7.} \label{ex:games:7} Consider the game for $\forall x.\; (P(x) \vee \lnot P(x))$ over a finite universe $\{a_1, \ldots, a_n\}$. Show that Prover has a winning strategy as long as Prover can decide, for each $a_i$, whether $P(a_i)$ holds. Now imagine $P$ is a property that requires checking infinitely many cases (like ``$x$ is the code of a Turing machine that halts on every input''). Why does the strategy break down? Connect this to the discussion of constructive vs.\ classical logic in the chapter.

\section*{Further reading}

If you want to follow the threads from games into theoretical computer science and modern cryptography:

\begin{itemize}
  \item Jaakko Hintikka, ``Language-Games for Quantifiers,'' \emph{American Philosophical Quarterly} Monograph Series no.\ 2 (1968), 46--72. Hintikka's own short and accessible exposition of how quantifiers are best read as a two-player game between a Verifier and a Falsifier---the philosophical source of the framework in this chapter.
  \item Shafi Goldwasser, Silvio Micali, and Charles Rackoff, ``\href{https://people.csail.mit.edu/silvio/Selected\%20Scientific\%20Papers/Proof\%20Systems/The_Knowledge_Complexity_Of_Interactive_Proof_Systems.pdf}{The Knowledge Complexity of Interactive Proof Systems},'' \emph{SIAM Journal on Computing} \textbf{18}(1), 186--208 (1989). The foundational paper on interactive proofs and zero-knowledge. The first few sections, where they set up the Prover/Verifier game, are surprisingly accessible and put a precise mathematical frame around the picture sketched in the chapter.
  \item Oded Goldreich, \emph{\href{https://www.wisdom.weizmann.ac.il/~oded/zk-tut02.html}{Zero-Knowledge: A Tutorial}}, available on the author's homepage at the Weizmann Institute. Written explicitly to be the most accessible on-ramp to zero-knowledge proofs for non-specialists. If interactive proofs caught your attention, this is where to go next.
\end{itemize}

\end{document}
